\documentclass[11pt]{article}
\usepackage[a4paper,margin=27mm]{geometry}
\usepackage{amsmath,amssymb,amsthm,xcolor,booktabs,array}
\newcommand{\PROVED}{\textcolor{green!40!black}{\textbf{PROVED}}}
\newcommand{\OPEN}{\textcolor{red!70!black}{\textbf{OPEN}}}
\newcommand{\AUDIT}{\textcolor{orange!70!black}{\textbf{AUDIT}}}
\newtheorem{theorem}{Theorem}
\newtheorem{proposition}{Proposition}
\title{Dual interpolation and the Abel--Poisson endpoint (v105)}
\date{13 August 2026}

\begin{document}
\maketitle

\section{The dual contractive endpoints}

On $L^p(\mathbb R_+,dx)$ let $U_nf(x)=f(nx)$ and
$S_n=\sqrt nU_n$.  Then
\begin{equation}
 \|U_n\|_{p\to p}=n^{-1/p},\qquad
 \|S_n\|_{p\to p}=n^{1/2-1/p}.                       \tag{1}
\end{equation}
The Banach adjoint of $S_n$ is
\begin{equation}
                  S_n'=n^{-1/2}U_{1/n}.              \tag{2}
\end{equation}
Thus $S_n$ is contractive for $1\le p\le2$, while its adjoint orientation
is contractive on the dual exponent $p'\ge2$.

For $\sigma>0$ define the Abel--Poisson operator
\begin{equation}
 P_\sigma:=\frac1{\zeta(1+\sigma)}
                \sum_{n\ge1}n^{-\sigma}U_n.          \tag{3}
\end{equation}

\begin{theorem}[Endpoint interpolation contraction --- \PROVED]
The operator $P_\sigma$ is a positive contraction on $L^1(dx)$, its Banach
adjoint is a positive contraction on $L^\infty(dx)$, and it is a contraction
on every interpolated $L^p$, in particular on $L^2(dx)$.
\end{theorem}

\begin{proof}
Since $\|U_n\|_{1\to1}=n^{-1}$,
\[
 \|P_\sigma\|_{1\to1}
 \le\frac{\sum_n n^{-1-\sigma}}{\zeta(1+\sigma)}=1.
\]
Positivity is termwise.  The Banach adjoint has the same norm, and
Riesz--Thorin interpolation proves the intermediate bounds.
\end{proof}

\section{The central limit}

The unitary logarithmic map
\[
 (Wf)(t)=e^{t/2}f(e^t):L^2(\mathbb R_+,dx)\longrightarrow L^2(\mathbb R,dt)
\]
sends $U_n$ to $n^{-1/2}$ times translation by $\log n$.  Hence the Fourier
multiplier of (3) is
\begin{equation}
 m_\sigma(\tau)=
 \frac{\zeta(\sigma+\tfrac12+i\tau)}{\zeta(1+\sigma)},          \tag{4}
\end{equation}
up to the harmless sign convention for $\tau$.

\begin{theorem}[The contractive Abel family collapses at the critical
endpoint --- \PROVED]
As $\sigma\downarrow0$, $P_\sigma\to0$ strongly on $L^2(dx)$.  Consequently
\begin{equation}
                 I-P_\sigma^*P_\sigma\longrightarrow I        \tag{5}
\end{equation}
strongly.  The limiting defect contains no odd Poisson quotient and no
row-(c) nuclear character.
\end{theorem}

\begin{proof}
For almost every real $\tau$, the numerator in (4) has a finite limit,
whereas $\zeta(1+\sigma)\sim\sigma^{-1}$.  Thus
$m_\sigma(\tau)\to0$ almost everywhere.  The preceding theorem gives
$|m_\sigma|\le1$ as an $L^2$ multiplier.  Dominated convergence in
Plancherel proves strong convergence to zero, and then (5) follows.
\end{proof}

The two primitive Tate conditions do not change this conclusion: they are
two evaluations of the analytic Mellin transform away from almost every
point on the real Plancherel axis.  They define a dense subspace of $L^2$
and $P_\sigma$ still converges strongly to zero there.

\section{Removing the normalization}

Multiplying (3) by $\zeta(1+\sigma)$ gives
\begin{equation}
 Z_\sigma=\sum_{n\ge1}n^{-\sigma}U_n,                 \tag{6}
\end{equation}
whose multiplier tends formally to
$\zeta(1/2+i\tau)$.  But
\begin{equation}
                  \|Z_\sigma\|_{L^1\to L^1}
                  =\zeta(1+\sigma)\longrightarrow\infty.      \tag{7}
\end{equation}
Thus the only rescaling which retains the Zeta boundary destroys the
endpoint contractivity used in the interpolation proof.

On critical $L^2$, the maximal multiplication operator by
$\zeta(1/2+i\tau)$ has dense range: its adjoint kernel consists of functions
supported on the real zero set, which has Lebesgue measure zero.  Therefore
its reduced Hilbert cokernel is zero.  The algebraic/non-Hausdorff cokernel
can retain evaluation jets, but is not a Hilbert quotient.

\begin{corollary}[Interpolation pivot verdict --- \PROVED]
The dual $L^1$--$L^\infty$ interpolation produces either
\begin{enumerate}
\item a contractive family whose critical limit is zero; or
\item the nontrivial Zeta multiplier with divergent endpoint norm and zero
      reduced Hilbert cokernel.
\end{enumerate}
It therefore does not construct the positive odd Hilbert completion with
the row-(c) trace.
\end{corollary}

\begin{center}
\begin{tabular}{>{\raggedright\arraybackslash}p{.58\linewidth}
                >{\centering\arraybackslash}p{.15\linewidth}
                >{\raggedright\arraybackslash}p{.17\linewidth}}
\toprule
Statement & Status & Location \\
\midrule
Dual endpoint contractions & \PROVED & (1)--(3) \\
Critical strong limit of normalized family & \PROVED & (4)--(5) \\
Trace-preserving odd defect & false for this family & Theorem 2 \\
Unnormalized endpoint contractivity & false & (7) \\
Contractivity (8c) on Loewy factors & \OPEN & not supplied by interpolation \\
\bottomrule
\end{tabular}
\end{center}

\end{document}
